<%= heading2 %>{Umgang mit TODOs}

<%- bexample %>
\textmarker{Markierter Text.}
<%- eexample %>

Bei \verb1\textmarker1 wird nur die Textfarbe geändert, da dies auch bei einigen Worten gut funktioniert.

<%- bexample %>
\textcomment{Markierter Text.}{Kommentar dazu.}
<%- eexample %>

<% if (isThesis && (todo == "pdfcomment") && (reallatexcompiler.startsWith("lualatex"))) { -%>
<%- bexample %>
\hl{In Gelb hervorgehoben.}
Provided indirectly by pdfcomment.sty (soulpos).
<%- eexample %>

<% } -%>
<%- bexample %>
\modified{Manuelle Markierung für Text, der seit der letzten Version geändert wurde.}
<%- eexample %>

<%- bexample %>
Das ist ein Text.
\change{FL1: Text angepasst}{Geänderter Text}.
<%- eexample %>

<%- bexample %>
Hier nur ein Kommentar\sidecomment{Kommentar}.
<%- eexample %>

<%- bexample %>
\todo{Hier muss noch kräftig Text produziert werden}
<%- eexample %>

Mehrere Autoren können ihre Änderungen und Kommentare mit dem Paket \href{https://ctan.org/pkg/phfcc}{phfcc} markieren; die Initialen erscheinen einmal pro Seite am Rand.
Pro Autor wird in der Präambel ein Befehl definiert (\verb|\phfMakeCommentingCommand|); \verb|\phfDisableCommentingCommands| blendet für die finale Version alles aus.

<%- bexample %>
Die Quantenmechanik entwickelte sich dank der Arbeiten von Schrödinger, Heisenberg und Dirac rasch,
die sich von früheren Ideen von Planck \FA{und Einstein} \FA[Einstein darf hier nicht fehlen.] inspirieren ließen.
\FA*{Dieser Satz soll entfernt werden.} \FA!{Das ist wichtig.}
<%- eexample %>
